\documentclass[11pt,a4paper]{article}
\usepackage[margin=24mm]{geometry}
\usepackage[T1]{fontenc}
\usepackage{lmodern,microtype}
\usepackage{amsmath,amssymb,amsthm,mathtools,booktabs}
\usepackage[hidelinks]{hyperref}
\hypersetup{pdftitle={A Jordan obstruction and generic irreducibility for Long--Moody braid representations},
pdfauthor={Haobo Ma and Wenlin Zhang}}
\setlength{\parskip}{0pt}
\setlength{\emergencystretch}{2em}
\newtheorem{theorem}{Theorem}
\newtheorem{proposition}[theorem]{Proposition}
\newtheorem{lemma}[theorem]{Lemma}
\newcommand{\C}{\mathbb C}
\newcommand{\Q}{\mathbb Q}
\newcommand{\Z}{\mathbb Z}
\newcommand{\F}{\mathbb F}
\newcommand{\Rcal}{\mathcal R}
\newcommand{\Tcal}{\mathcal T}
\newcommand{\End}{\operatorname{End}}
\newcommand{\im}{\operatorname{im}}
\newcommand{\tr}{\operatorname{tr}}
\newcommand{\rank}{\operatorname{rank}}
\newcommand{\Bur}{\mathrm{Bur}^{\mathrm{red}}}
\title{A Jordan obstruction and generic irreducibility\\
for Long--Moody braid representations}
\author{Haobo Ma\thanks{ChronoAI Pte Ltd; The Omega Institute.}
\and Wenlin Zhang\thanks{National University of Singapore; The Omega Institute.}}
\date{10 October 2026}
\begin{document}
\maketitle
\begin{abstract}
We resolve the Jones--Temperley--Lieb Jordan-obstruction clause of
Negami's Problem~6.8 and extend its generic irreducibility result to
$n=5$. For every untwisted endpoint-sector seed of dimension at least
two, with $n\ge3$ and $\ell\ge4$, an explicit four-dimensional block
contains a nontrivial length-two Jordan chain. The chain survives every
admissible middle-convolution quotient, so both braid images are
unbounded and admit no positive-definite invariant Hermitian form.
For the separate character-twisted Burau family, an exact minor of order
$893$ proves generic irreducibility of the twenty-dimensional summand
$\Rcal_5$ in the positive parameter region. We construct an invariant
$10+10$ splitting along the entire admissible diagonal $\tau=q$ and
prove that every ambient minor of order $893$ is divisible by
$(\tau-q)^3$. Exact cyclotomic certificates establish irreducibility at
$95$ off-diagonal points and identify the two irreducible constituents
at seven diagonal points.
\end{abstract}

\section{The construction and the two questions}

Negami constructs braid representations from character-twisted Burau
seeds and asks in \cite[Problem~6.8]{Negami} for the remaining
irreducibility locus, the image closure on irreducible sectors, and
an exact Jordan obstruction for the Jones--Temperley--Lieb seeds of
\cite[Remark~6.7(b)]{Negami}. We prove the Jordan obstruction uniformly
in the endpoint sector, including its survival on the convolution
quotient. For the Burau family we extend
\cite[Proposition~5.10]{Negami} from $n=3,4$ to $n=5$ and determine an
explicit reducibility curve. These are two different seed families.

We use the conventions of \cite[Definition~1.3 and Definitions~1.4--1.7]{Negami}.
A seed consists of invertible operators $g_1,\ldots,g_n$ and
$s_1,\ldots,s_{n-1}$ on a vector space $V$ satisfying the seed
relations (the braid and Artin compatibility relations of Definition~1.3).
The ambient
Long--Moody module is $U=V^{\oplus n}$. Its braid generator $S_i$
acts by $s_i$ on every slot other than $i,i+1$, and
\begin{equation}\label{eq:S}
 (S_iv)_i=s_ig_iv_{i+1},\qquad
 (S_iv)_{i+1}=s_i\bigl(v_i+(I-g_{i+1})v_{i+1}\bigr).
\end{equation}
The middle-convolution module, for $\lambda\in\C\setminus\{0,1\}$, is
\[
 Q_\lambda=U/(K_{\rm seed}+L_\lambda),\qquad
 K_{\rm seed}=\bigoplus_{j=1}^n\ker(g_j-I),
\]
where, with $P_n=g_1\cdots g_n$,
\begin{equation}\label{eq:L}
 L_\lambda=\{z:z_j=g_{j+1}z_{j+1}\ (1\le j<n),\quad
                    z_n=\lambda P_nz_n\}.
\end{equation}
These subspaces are braid invariant. The operators in \eqref{eq:S}
are independent of $\lambda$ \cite[Proposition~2.3]{Negami}.

\section{The endpoint-sector Jordan obstruction}

Fix integers $n\ge3$ and $\ell\ge4$. An endpoint sector $V$ has an
orthonormal basis indexed by paths
$p=(p_0,\ldots,p_{n+1})$ on $A_{\ell-1}$, with vertices
$1,\ldots,\ell-1$, $p_0=1$, and a fixed terminal vertex.
Consecutive vertices differ by one. Set
\[
 \theta=\pi/\ell,\qquad \mu_c=\sin(c\theta),\qquad
 \delta=2\cos\theta,\qquad
 \alpha=i e^{-i\theta/2},\qquad b=-\alpha^{-3}.
\]
The operator $E_j$, $0\le j<n$, changes vertex $j+1$ only.
Its entry from a path $r$ to a path $p$ is zero unless the paths
agree away from $j+1$ and $p_j=p_{j+2}$; otherwise
\begin{equation}\label{eq:E}
 (E_j)_{p,r}=
 \frac{\sqrt{\mu_{p_{j+1}}\mu_{r_{j+1}}}}{\mu_{p_j}}.
\end{equation}
The untwisted normalization of \cite[Remark~6.7(b)]{Negami} is
\[
 \rho_j=\alpha I+\alpha^{-1}E_j,\qquad s_j=\rho_j,
 \qquad g_1=\rho_0^2,\qquad
 g_{j+1}=\rho_jg_j\rho_j^{-1}.
\]

\begin{theorem}\label{thm:jordan}
For every endpoint sector $V$ of dimension at least two,
$K_{\rm seed}=0$, and there are $w,v\in U$ such that
\[
 (S_1-bI)w=v\ne0,\qquad (S_1-bI)v=0.
\]
The images of $w,v$ give the same nontrivial Jordan chain on
$Q_\lambda$ for every $\lambda\in\C\setminus\{0,1\}$.
Both the ambient and quotient braid images are unbounded and preserve
no positive-definite Hermitian form.
\end{theorem}

\begin{proof}
Put $\mu_0=\mu_\ell=0$. The sine addition formula gives
$\mu_{c-1}+\mu_{c+1}=\delta\mu_c$, also at the boundary vertices.
Every nonzero local block of $E_j$ is $uu^{\mathsf T}/\mu_c$,
where $u_d=\sqrt{\mu_d}$ for neighbors $d$ of $c$. Thus
$E_j^2=\delta E_j$. Since
$\delta=-(\alpha^2+\alpha^{-2})$, we obtain
\[
 (\rho_j-\alpha I)(\rho_j+\alpha^{-3}I)=0,
 \qquad (g_j-\alpha^2I)(g_j-\alpha^{-6}I)=0.
\]
The second identity follows first for $g_1$ by squaring, then for
all $g_j$ by conjugation. The seed operators are invertible because
the two roots $\alpha,-\alpha^{-3}$ are nonzero.
Now $\alpha^2=-e^{-i\theta}$ and $\alpha^6=-e^{-3i\theta}$ have
positive imaginary parts for $0<\theta\le\pi/4$.
Consequently neither $\alpha^2$ nor $\alpha^{-6}$ is one.
Applying the polynomial identity to a fixed vector of $g_j$ shows
that it is zero, proving $K_{\rm seed}=0$.

There is at most one strictly ascending path. Since the sector has
at least two paths, some path has a downward step. At the first such
step, delete the two-edge detour $c-1,c,c-1$ and prepend the loop
$1,2,1$. This preserves length and endpoint and produces a path
beginning $1,2,1,2$. Here the fourth vertex is present because
$n+1\ge4$, and a step after vertex $1$ must go to $2$.
Replacing the third vertex by $3$ gives a second legal path beginning
$1,2,3,2$, since $\ell\ge4$. Denote these paths by $p,p'$;
their tails agree. On their span,
\[
 E_0=\begin{pmatrix}\delta&0\\0&0\end{pmatrix},\qquad
 E_1=\frac1{\mu_2}
 \begin{pmatrix}\mu_1&\sqrt{\mu_1\mu_3}\\
 \sqrt{\mu_1\mu_3}&\mu_3\end{pmatrix}.
\]
Put $f_0=e_p$ and $f_1=\alpha^2E_1f_0$.
The coefficient of $e_{p'}$ in $f_1$ is nonzero, so $f_0,f_1$
are independent. Using $\mu_2=\delta\mu_1$ gives
$E_0f_1=\alpha^2f_0$ and $E_1f_1=\delta f_1$.
In this basis the first two seed matrices are
\begin{equation}\label{eq:RU}
 R=\rho_0=\begin{pmatrix}-\alpha^{-3}&\alpha\\0&\alpha\end{pmatrix},
 \qquad U=\rho_1=\begin{pmatrix}\alpha&0\\\alpha^{-3}&-\alpha^{-3}\end{pmatrix},
 \qquad U^{-1}=\begin{pmatrix}\alpha^{-1}&0\\\alpha^{-1}&-\alpha^3\end{pmatrix}.
\end{equation}
Let $M$ be the four-dimensional space consisting of this span in
slots $1,2$ and zero in all other slots. It is $S_1$ invariant, with
\begin{equation}\label{eq:T}
 T=S_1|_M=\begin{pmatrix}0&UG\\U&U(I-H)\end{pmatrix},
 \qquad G=R^2,\qquad H=UGU^{-1}.
\end{equation}
In the order $(f_0,f_1)$ in slot $1$, then $(f_0,f_1)$ in slot $2$,
put $d=\alpha^6-1$ and set
\begin{equation}\label{eq:wv}
 w=\begin{pmatrix}
 -\alpha^5(\alpha^2+1)(\alpha^8-\alpha^6-\alpha^4-\alpha^2+1)\\
 \alpha(2\alpha^6+\alpha^4-1)\\
 -\alpha^{11}(\alpha^2+1)\\0
 \end{pmatrix},\qquad
 v=\begin{pmatrix}-\alpha^2(\alpha^4-1)d\\\alpha^{-2}d\\0\\d\end{pmatrix}.
\end{equation}
Multiplying the matrices in \eqref{eq:RU}--\eqref{eq:T} and
collecting Laurent powers gives the identities
\begin{equation}\label{eq:chain}
 (T+\alpha^{-3}I)w=v,\qquad (T+\alpha^{-3}I)v=0
\end{equation}
for every nonzero $\alpha$. At the specified phase, $d\ne0$,
so the last coordinate of $v$ is nonzero. Extension by zero to the
remaining slots proves the ambient chain.

If $z\in M\cap L_\lambda$, then $z_3=\cdots=z_n=0$.
The recurrence in \eqref{eq:L} gives $z_2=g_3z_3=0$ and
$z_1=g_2z_2=0$. Hence $M\cap L_\lambda=0$.
Since $K_{\rm seed}=0$, the quotient map is injective on $M$,
and \eqref{eq:chain} descends with $[v]\ne0$.
This argument includes the parameters for which $L_\lambda\ne0$.

Finally $|b|=1$, and for every integer $k\ge1$ the chain gives
$S_1^kw=b^kw+kb^{k-1}v$. In any norm,
\[
 \|S_1^kw\|\ge k\|v\|-\|w\|\longrightarrow\infty.
\]
The same inequality holds on the quotient. An invariant positive-definite
Hermitian form would make every power an isometry, a contradiction.
\end{proof}

All standard braid generators are conjugate, so each has the same
Jordan obstruction. A
\href{https://github.com/the-omega-institute/trureturing/blob/868ae8ef94f515b39323d972ccfd7120036bff79/D5/S3/Quantum/Algebra/JonesTemperleyLiebLongMoodyJordanObstruction.lean}{formalization of the ambient assertion}
covers $K_{\rm seed}=0$ and the existence of the length-two chain with
the path conventions above. The quotient and unboundedness conclusions
follow from the written argument.

\section{The character-twisted Burau summand}

For the remainder of the paper set
\[
 \Tcal_n=\{(a,c):a>0,\ c>0,\ nc+(n+1)a<1\},\qquad
 q=e^{2\pi ia},\qquad \tau=e^{2\pi ic}.
\]
Let $B_j(q)$, $0\le j<n$, be the unreduced Burau matrix on
$\C^{n+1}$ with local block
$\left(\begin{smallmatrix}1-q&q\\1&0\end{smallmatrix}\right)$
at coordinates $j,j+1$ and identity elsewhere. Define
\[
 s_i=B_i\ (1\le i<n),\qquad g_1=\tau B_0^2,
 \qquad g_{j+1}=B_jg_jB_j^{-1}.
\]
Write $U_0=(\C^{n+1})^{\oplus n}$ for the resulting ambient module.
The reduced seed is $V=\ker\ell_q$, where
$\ell_q(x)=\sum_{j=0}^n q^jx_j$; its integral Laurent basis is
$b_j=e_j-q^je_0$, $1\le j\le n$. On $V^{\oplus n}$ put
\[
 D(v)=\sum_{j=1}^n(g_j-I)v_j,\qquad \Rcal_n=\ker D.
\]
On $\Tcal_n$, $g_1-I$ is invertible, so $D$ is surjective and
$\dim\Rcal_n=n(n-1)$. By \cite[Proposition~5.9]{Negami},
the map $(\Psi v)_j=v_j-g_{j+1}v_{j+1}$ identifies $\Rcal_n$
with $V^{\oplus(n-1)}$. Indeed, its kernel is $F(V)$, where
\[
 F(v)=(g_2\cdots g_nv,\ldots,g_nv,v),\qquad
 DF=P_n-I,
\]
and $P_n-I$ is invertible on $\Tcal_n$
\cite[Theorem~5.3 and Remark~5.8]{Negami}.
In these coordinates the generator is
\begin{equation}\label{eq:A}
 (A_iu)_j=
 \begin{cases}
 s_i(u_{i-1}+g_iu_i),&j=i-1,\\
 -s_ig_iu_i,&j=i,\\
 s_i(u_i+u_{i+1}),&j=i+1,\\
 s_iu_j,&\text{otherwise},
 \end{cases}
\end{equation}
with nonexistent endpoint rows omitted and all seed operators restricted
to $V$. For $n=5$, the ambient and reduced dimensions are $30$ and $20$.

\section{A nonzero Laurent minor and generic irreducibility}

For row-major vectorization of a $d\times d$ matrix $X$, the
commutant system for generators $M_1,\ldots,M_{n-1}$ is
\begin{equation}\label{eq:C}
 \mathcal C(M)=\begin{pmatrix}
 M_1\otimes I_d-I_d\otimes M_1^{\mathsf T}\\
 \vdots\\
 M_{n-1}\otimes I_d-I_d\otimes M_{n-1}^{\mathsf T}
 \end{pmatrix}.
\end{equation}
Its kernel is precisely $\End_{B_n}(U)$.
Write $\mathcal C_5=\mathcal C(S)$ for the $3600\times900$
ambient system and $\mathcal C^{\rm red}_5=\mathcal C(A)$ for the
$1600\times400$ reduced system. Their entries belong to
$\Z[q,q^{-1},\tau]$.

We fix the minors by the following deterministic selection rule.
Evaluate at $(q,\tau)=(2,3)$ over $\F_{1009}$ and perform forward
elimination, scanning columns from left to right. In each pivot column,
choose the first nonzero row at or below the current pivot position,
swap it into that position, normalize its pivot to one, and eliminate
entries below. Record the original labels of the pivot rows and the
pivot columns in their selection order. All initial labels are zero
based, with rows ordered by generator and then by matrix entry.
Let $f_5$ and $g_5$ be the determinants of the selected submatrices of
the Laurent systems $\mathcal C_5$ and $\mathcal C^{\rm red}_5$,
respectively. The certificate repository at
\url{https://github.com/the-omega-institute/jtl-seeds-problem68} gives
their full sparse Laurent matrices as well as the index lists.

\begin{theorem}\label{thm:generic}
The Laurent polynomial $f_5$ is nonzero. For every
$(a,c)\in\Tcal_5$ such that
$f_5(e^{2\pi ia},e^{2\pi ic})\ne0$, the representation $\Rcal_5$
is irreducible over $\C$. This nonvanishing set is dense and open in
$\Tcal_5$. The same conclusions hold with $g_5$ in place of $f_5$.
\end{theorem}

\begin{proof}
Exact elimination at $(2,3)$ over $\F_{1009}$ gives
$\rank\mathcal C_5=893$ and the selected determinant $992$.
Thus $f_5$ cannot be the zero integer Laurent polynomial.
Its nonvanishing bounds the ambient commutant dimension above by
$900-893=7$.
On $\Tcal_5$, \cite[Proposition~5.7]{Negami} gives
\begin{equation}\label{eq:decomp}
 U_0\cong\mathbf1^{\oplus2}\oplus\Bur_5(\tau)
       \oplus\Bur_5(q)\oplus\Rcal_5(q,\tau).
\end{equation}
Consequently
\[
 \dim\End_{B_5}(U_0)\ge4+1+1+\dim\End_{B_5}(\Rcal_5)\ge7.
\]
Equality forces $\End_{B_5}(\Rcal_5)=\C I$.
For $(a,c)\in\Tcal_5$, choose $l=(1-5c-6a)/2$ and
$\lambda=e^{2\pi i l}$. Then $0<l<1-5c-6a<1$, so Negami's Theorem~5.3 supplies a
positive-definite invariant form on the full $30$-dimensional module,
and Proposition~5.4(1) gives unitarizability without an irreducibility
hypothesis. The braid matrices in \eqref{eq:S} are independent of this
choice of $\lambda$.
Its restriction to $\Rcal_5$ is positive definite. Any invariant
subspace therefore has an invariant orthogonal complement. A proper
nonzero invariant subspace would give a nonscalar commuting projection,
proving the claimed irreducibility.

For the reduced system the exact rank is $399$ and the selected
determinant is $742$ modulo $1009$. Thus $g_5\ne0$; its nonvanishing
forces a scalar reduced commutant directly, and the same positive-form
argument applies. The rational witness $(2,3)$ is used to prove
nonvanishing, while the irreducibility conclusion is made on $\Tcal_5$.

A nonzero Laurent polynomial cannot vanish on an open patch of the
parameter torus. Its restriction is a real analytic finite Fourier
sum; vanishing on a patch would force vanishing on the torus and then
all coefficients to be zero. Hence each nonvanishing set is dense and
open in $\Tcal_5$.
\end{proof}

For comparison with \cite[Proposition~5.10]{Negami}, the table gives
the exact ranks and nullities at $(2,3)$ over $\F_{1009}$ and the
selected determinant residues there and at $(5,7)$ over $\F_{1019}$.
\begin{center}
\begin{tabular}{@{}lrrrrr@{}}
\toprule
Module & Dimension & Rank & Nullity & $\F_{1009}$ & $\F_{1019}$\\
\midrule
$U_0$, $n=3$ & 12 & 137 & 7 & 786 & 18\\
$U_0$, $n=4$ & 20 & 393 & 7 & 962 & 38\\
$U_0$, $n=5$ & 30 & 893 & 7 & 992 & 393\\
$\Rcal_5$ & 20 & 399 & 1 & 742 & 471\\
\bottomrule
\end{tabular}
\end{center}
The $n=3,4$ rows reproduce the earlier certificate values; the $n=5$
rows supply the new generic irreducibility certificates.

\section{The diagonal splitting and factors of the minors}

Write $[m]_q=1+q+\cdots+q^{m-1}$.
An exact polynomial commuting operator gives the following splitting.

\begin{lemma}\label{lem:K}
There is a matrix $K(q)\in\operatorname{Mat}_{20}(\Z[q])$,
of entry degree at most eleven, such that
\begin{equation}\label{eq:K}
 [K,A_i(q,q)]=0\ (1\le i\le4),\qquad
 K^2=q[6]_q^2I,\qquad \tr K=0,
\end{equation}
and, in zero-based indices, $K_{1,0}=q[3]_q$.
\end{lemma}

\begin{proof}
Here is a finite specification of $K$ using the matrices in
\eqref{eq:A}. For each integer $v=2,\ldots,13$, solve the rational
commutant equations for $A_i(v,v)$, together with
$\tr K_v=0$ and $(K_v)_{1,0}=v[3]_v$.
These systems have unique solutions. Define
\begin{equation}\label{eq:interp}
 K(q)=\sum_{v=2}^{13}K_v
       \prod_{\substack{2\le w\le13\\w\ne v}}\frac{q-w}{v-w}.
\end{equation}
The resulting entries have integer coefficients. Expansion of the
four commutators, the square residual, the trace, and the specified
entry verifies \eqref{eq:K} coefficient by coefficient over
$\Z[q,q^{-1}]$. The certificate records every entry of $K$.
For completeness, the reduced commutant rank at each of the twelve
interpolation points is $398$ modulo $1019$. This gives rational
commutant dimension at most two; $I$ and the verified $K(v)$ give
the reverse bound. Since $K(v)$ has trace zero and prescribed nonzero
entry, it is the unique normalized solution used in \eqref{eq:interp}.
Thus interpolation specifies the candidate, and the polynomial
identities establish the assertion for all $q$.
\end{proof}

\begin{theorem}\label{thm:diagonal}
For every $0<a<1/11$, at $q=\tau=e^{2\pi ia}$,
\[
 \Rcal_5=W_+\oplus W_-,\qquad \dim W_+=\dim W_-=10.
\]
These invariant subspaces are the images, in the coordinates of
\eqref{eq:A}, of
\begin{equation}\label{eq:projectors}
 E_\pm=\frac12\left(I\pm\frac{K(q)}{q^{1/2}[6]_q}\right),
 \qquad q^{1/2}=e^{\pi ia}.
\end{equation}
\end{theorem}

\begin{proof}
The admissible diagonal in $\Tcal_5$ is exactly $0<a<1/11$.
Here $q\ne1$ and $0<6a<1$, so $[6]_q=(q^6-1)/(q-1)\ne0$.
Equation~\eqref{eq:K} gives complementary idempotents:
$E_\pm^2=E_\pm$, $E_+E_-=0$, and $E_++E_-=I$.
They commute with all generators. Their traces, hence their ranks,
are ten. Their images are the required invariant summands.
The isomorphism $\Psi|_{\ker D}$ transfers them to the augmentation
kernel.
\end{proof}

\begin{proposition}\label{prop:factor}
Every minor of order $893$ of $\mathcal C_5$ is divisible by
$(\tau-q)^3$ in $\Z[q,q^{-1},\tau]$. Every minor of order $399$
of $\mathcal C^{\rm red}_5$ is divisible by $\tau-q$.
\end{proposition}

\begin{proof}
On the admissible diagonal the two Burau summands in
\eqref{eq:decomp} are isomorphic, and scalar maps between the two
copies give four independent endomorphisms. The two trivial copies
give four more, and $E_+,E_-$ give two. Hence the ambient commutant
has dimension at least ten and the reduced commutant at least two.
Their ranks are at most $890$ and $398$, respectively.
The admissible arc is Zariski dense in the curve $\tau=q$, so these
rank bounds hold over $\Q(q)$ after setting $\tau=q$.

Localize $\Q(q)[\tau]$ at $\tau-q$. A square submatrix of order
$893$ whose reduction has rank at most $890$ has at least three
invariant factors divisible by the uniformizer $\tau-q$. Its
determinant is therefore divisible by the cube. The reduced argument
uses orders $399$ and $398$. Finally, division by the monic polynomial
$(\tau-q)^3$ over $\Z[q,q^{-1}]$ has a remainder of degree less
than three in $\tau$. Its image in $\Q(q)[\tau]$ is zero, so the
remainder is zero. Division by $\tau-q$ gives the reduced assertion.
\end{proof}

\section{Exact cyclotomic points and the remaining questions}

Consider all $102$ points
\[
 (a,c)=(j/84,k/84),\qquad j,k>0,\qquad 6j+5k<84.
\]
Set $\zeta=e^{2\pi i/84}$. The map $\zeta\mapsto527\in\F_{1009}$
defines a homomorphism $\Z[\zeta]\to\F_{1009}$, since
$\Phi_{84}(527)=0$ modulo $1009$.
The element $527$ has order $84$: its $84$th power is one, while
its powers $84/r$ for $r=2,3,7$ are not one.
Thus a nonzero reduced minor after this map certifies a nonzero
cyclotomic minor in characteristic zero.

At all $95$ off-diagonal points, $g_5$ has nonzero image modulo
$1009$, proving irreducibility over $\C$ by
Theorem~\ref{thm:generic}. At the seven diagonal points
$j=k=1,\ldots,7$, the reduced modular rank is $398$.
The modular rank bounds the complex commutant dimension above by two;
$I$ and $K$ give two independent complex endomorphisms.
Its dimension is therefore exactly two. Complete reducibility makes
this dimension the sum of squares of irreducible multiplicities.
Together with the invariant $10+10$ splitting, this proves that
$W_+$ and $W_-$ are irreducible and inequivalent at these seven points.
At $j=k=1$ the ambient modular rank is $890$, so its generic rank
along the diagonal is exactly $890$.

The verifier selects an order-$398$ minor of the full reduced commutant
system at $q=\tau=2$ over $\F_{1019}$, restricts it to $\tau=q$, and
checks its determinant at $q=2$ exactly over $\Z$. This gives a nonzero
univariate Laurent polynomial $h(q)$. Whenever $h(q)\ne0$, the reduced
commutant has dimension two by the same $I,K$ argument, so complete
reducibility gives the corresponding conclusion for $W_+$ and $W_-.$

The exact finite computations supporting the paper are the commutant
ranks and selected determinant residues, the Laurent coefficient
identities in Lemma~\ref{lem:K} and \eqref{eq:chain}, and the
cyclotomic reductions at the stated points. The Jordan obstruction
is uniform in $n,\ell$ and the sector; it follows from the local
construction, without a finite sector scan. For the Burau family,
the certificates prove irreducibility on
$\{f_5\ne0\}\cap\Tcal_5$ and $\{g_5\ne0\}\cap\Tcal_5$,
and the polynomial projectors prove reducibility along the entire
admissible diagonal.

The complete exceptional locus away from the diagonal, the residual
factors and their exact multiplicities, and the image-closure and
$\mathrm{SU}(20)$ questions in \cite[Problem~6.8]{Negami} remain open.
A zero of one selected minor alone does not imply reducibility. Since
$h(q)$ has only finitely many roots, $W_+$ and $W_-$ are irreducible
and inequivalent at all but finitely many admissible diagonal points;
the complete exceptional diagonal set remains undetermined.

\section*{Data availability}

The exact certificates and reproduction programs are available at
\url{https://github.com/the-omega-institute/jtl-seeds-problem68}.

\section*{Use of AI tools}

We used an AI orchestration system built on OpenAI's ChatGPT and Codex
and Anthropic's Claude. AI agents searched for proofs and counterexamples,
ran computations, and produced the Lean formalization of the ambient
Jordan obstruction. Separate ChatGPT and Codex runs double-checked
the arguments and computations; Claude coordinated the work.
AI tools also assisted in preparing this manuscript.

\begin{thebibliography}{1}
\bibitem{Negami}
Haru Negami,
\emph{Quantum gates from the middle convolution of twisted Burau
representations}, arXiv:2610.00293v1 (2026).
\url{https://arxiv.org/abs/2610.00293v1}.
\end{thebibliography}
\end{document}
